\section{Trait-Conditioned C--K Trajectory Modeling}
\label{sec:trait_conditioned_cktm}

We model individualized designer role-playing as the joint generation of an
observable C--K trajectory and a first-person decision with its justification.
The trajectory summarizes operations expressed in the response; it is not an
account of the model's hidden reasoning process.

Given a public design situation $X^\ast$, generation is conditioned on
designer $i$'s professional background $P_i$, Q2 reasoning record $H_i$,
situation-adaptive trait cues $\widetilde V_i(X^\ast)$, selected
trait-conditioned rules $R_i$ learned from Q1--Q6, and aligned examples $E_i$:
\begin{equation}
p_\psi(M_i,Y_i\mid
X^\ast,P_i,\widetilde V_i(X^\ast),H_i,R_i,E_i).
\label{eq:tcktm_joint}
\end{equation}
The complete 30-dimensional profile $V_i$ is used internally for cue and rule
selection, but is not fully serialized in the prompt. Instead,
$\widetilde V_i(X^\ast)$ contains the four dimensions that are most relevant
to Q7 and salient for the target designer. $E_i$ contains three
state--text-aligned Q7 examples from other designers; they calibrate operation
semantics and output granularity rather than provide answers to imitate. The
target designer is excluded from $E_i$, and their Q7 response is never
provided. Public task facts take precedence over all conditioning signals.

\subsection{Learning Trait-Conditioned Rules}

We derive trait-conditioned rules from Q1--Q6 and freeze them before Q7
generation. The rules express preferences over initial orientations,
subsequent operations, and situation-specific choices; none specifies a
deterministic transition. Initial-orientation rules apply to Seg1, whereas the
Seg2 method uses the remaining three families in
Table~\ref{tab:trait_conditioned_rules}.

\begin{table}[t]
\centering
\small
\caption{Trait-conditioned rule families.}
\label{tab:trait_conditioned_rules}
\begin{tabular}{p{0.22\columnwidth} p{0.36\columnwidth} p{0.32\columnwidth}}
\toprule
\textbf{Rule type} & \textbf{Condition} & \textbf{Guidance} \\
\midrule
Initial orientation & Trait tendency associated with the initial response. &
Prefer an initial C--K operation. \\
Single-trait transition & Trait tendency associated with a transition. &
Prefer a next operation given the current operation. \\
Trait-relation transition & Co-activation, tension, dominance, or balance
between two traits. & Adjust the next-operation preference beyond separate
trait effects. \\
Contextual preference & Trait tendency under a recurring design conflict. &
Prefer an operation appropriate to that conflict. \\
\bottomrule
\end{tabular}
\end{table}

For a Q2 trajectory, let $z_{it}$ denote designer $i$'s C--K operation at
position $t$. For each source--target pair, we fit an L1-regularized logistic
regression over the standardized V30 profile and relative trajectory position:
\begin{equation}
z_{i,t+1}\sim p(\,\cdot\mid z_{it},V_i,\phi_{it}).
\label{eq:tcktm_transition}
\end{equation}
Sparse coefficients identify trait tendencies associated with the next
operation after accounting for the current operation and position. Candidate
rules are filtered using participant-level support, bootstrap stability, sign
consistency, and held-out evidence.

Trait-relation rules are learned in two stages. We first identify stable trait
pairs using participant-level rank correlations. For each pair, we construct
interpretable predicates---co-activation, tension, dominance, and balance---
and test whether the predicate adds predictive information beyond the two
individual trait effects. Q3--Q6 provide contextual operation choices and are
used to learn contextual preferences rather than artificial transitions.
These rules are empirical soft preferences, not causal or universal
psychological laws. Detailed thresholds and rule statistics are provided in
Appendix~\ref{app:trait_conditioned_details}.

\subsection{One-Call Conditioned Joint Generation}

For designer $i$, the selector uses the complete $V_i$ internally and retains
at most one active Seg2 rule from each family: a single-trait transition rule,
a trait-relation transition rule, and a contextual preference rule. If
single-trait and relation rules conflict for the same source operation, the
relation rule receives soft priority. Rules may be ignored when inapplicable,
and public facts always override rule guidance.

The model produces all state--text pairs in one call:
\begin{equation}
\mathcal U_i=[(m_{i1},y_{i1}),\ldots,(m_{iT_i},y_{iT_i})],
\quad
M_i=[m_{i1},\ldots,m_{iT_i}],
\quad
Y_i=y_{i1}\oplus\cdots\oplus y_{iT_i}.
\end{equation}
Each $y_{it}$ must realize the observable operation $m_{it}$. In particular,
$C\!\rightarrow\!K$ formulates a future or conditional verification step
without inventing evidence, whereas $K\!\rightarrow\!K$ derives consequences
from known facts without introducing a new proposal. The concatenated
$Y_i$ is a coherent first-person Chinese response and does not expose trait
scores or rule instructions.
